\paragraph{LZ Detector and Calibrations --}
The heart of the LZ experiment is a low-background dual-phase time projection chamber (TPC) containing 7 tonnes of active liquid xenon (LXe)~\cite{LZ:Experiment_2020, LZ:TDR_2017}.
The TPC is equipped with arrays of photomultiplier tubes (PMTs) at its top and bottom to detect optical signals.
Cathode, gate, and anode electrodes create a set of electric fields that drift ionization electrons to the liquid surface and extract them into the gas phase.
The liquid-gas interface lies between the gate and anode at the top of the detector, and the cathode defines the bottom of the active region. A bottom electrode sits below the cathode to shield the lower PMT array.

The TPC  is surrounded by a LXe Skin veto detector, consisting of 2~tonnes of LXe instrumented for scintillation readout whose main purpose is to identify $\gamma$-rays  entering or exiting the TPC.
The TPC and Skin are contained within a vacuum-insulated cryostat, which is in turn surrounded by the Outer Detector (OD) veto. The OD is 17~tonnes of gadolinium-loaded liquid scintillator (0.1\% by mass) and 229~tonnes of ultra-pure water, read out by PMTs in the water space, and it detects neutrons, $\gamma$-rays, and muons.
The LZ detectors are described in more detail in Ref.~\cite{LZ:Experiment_2020}.

Energy depositions in the TPC typically produce two signals: prompt scintillation light (S1), produced at the location of the interaction; and delayed electroluminescence light (S2), produced when extracted electrons are accelerated through the gaseous phase by the applied electric field.
The S2 light hit pattern in the top PMT array enables position reconstruction in the ($x,y$) plane, and the time interval between S1 and S2 pulses provides the depth, $z$. S1 and S2 signal sizes are measured in photons detected (phd) and vary with the position of the interaction vertex. Position-corrected values, denoted S1$c$ and S2$c$, allow for scale factors to correlate the signals to the original number of photons and electrons produced.
The photon gain factor $g_1=0.110\pm0.002$~phd/photon and the electron gain factor $g_2=34.5\pm1.1$~phd/electron are determined using monoenergetic peaks from background and calibration sources.

The S1 to S2 ratio provides discrimination power between electron recoil (ER) events, produced mainly by $\gamma$-ray and $\beta$-decay interactions, and nuclear recoil (NR) events, expected from dark matter interactions as well as neutrons and coherent scattering of neutrinos -- as illustrated in~\autoref{fig:calibrations} with calibration datasets. 
The ER response is measured using \textit{in-situ} $\beta$-decays of tritium and $^{14}$C, injected as methane, and $^{212}$Pb from $^{220}$Rn injections.
The NR response is measured using neutrons from Deuterium-Deuterium (D-D) fusion (0--74~keV NRs) and from AmBe ($\alpha$,n) reactions (0--330~keV NRs)~\cite{LZ:Calibrations_2024}. 
We define the ``ER band'' and ``NR band'' as the distributions in \{S1$c$, log$_{10}$(S2$c$)\} space populated by flat-in-energy ER and NR recoil spectra, shown by the solid and dashed lines in Fig.~\ref{fig:calibrations}. The detector and model response parameters from NEST v2.4.5 (Noble Element Simulation Technique)~\cite{ NEST:paper_2023, nest_v2_4_5} are tuned to describe the ER and NR calibration data. 
The best-fit NEST parameters are provided in the \textcolor{black}{Data Release}, with more detail of the model in the \textcolor{black}{Supplemental Material}. 







 


\begin{figure}
    \centering
    \includegraphics[width=0.95\linewidth]{\plotfolder Fig2_Calibrations.pdf}
    \caption{
    Calibration events in \{S1$c$,~log$_{10}$(S2$c$)\} space from tritium and $^{14}$C (light blue), $^{212}$Pb (green), D-D neutrons (orange), and AmBe neutrons (purple). 
    The mean (solid) and 10\%-90\% percentiles (dashed) of the ER (dark blue) and NR (red) bands are given by the calibrated NEST model, taking flat energy spectra and evaluating along lines of fixed S$1c$. The  solid grey lines show contours of constant energy in \{S1$c$, log$_{10}$(S2$c$)\} space. The dashed grey line indicates the WS~ROI.
    Tritium and ${}^{14}$C data are cut off above S1$c=500$~phd to avoid ${}^{133}$Xe contamination from activation due to a preceding neutron calibration. 
    To increase statistics, the calibration datasets include events outside the fiducial volume used in the main WIMP search analysis, leading to some wall leakage. Grey points are outside of $4\sigma$ from the mean of the ER and NR bands. These points are not used in deriving the NEST model as they are likely from interactions near the TPC walls.  
    }
    \label{fig:calibrations}
\end{figure}